\begin{description}
\item[(T05.1)] Angular diameter distance $= D_L/(1+z)^2$
  \begin{align*}
  \text{Scale} &= \text{Angular diameter distance} \times
    \frac{\pi}{60 \times 60 \times 180}~\mathrm{ly/arcsecond} \\
  &= 20.7~\mathrm{ly~milliarcsecond^{-1}}.
  \tag*{(1.0)}
  \end{align*}
  \begin{center}
  \begin{tabular}{cccc}
  Date of observation & $\phi_{\mathrm{blob}}$ (milliarcsecond)
    & $l_{\mathrm{blob}}$ (ly) & $\beta_{\mathrm{app}}$ \\
  \hline
  $T_0$            & 1.6 & 33 & $\times$ \\
  $T_0 + 605$~d    & 2.1 & 44 & 6.6 \\
  $T_0 + 1186$~d   & 2.3 & 48 & 2.5 \\
  $T_0 + 1786$~d   & 2.8 & 58 & 6.1 \\
  $T_0 + 2338$~d   & 3.1 & 64 & 4.0 \\
  \end{tabular}
  \hfill(3.0)
  \end{center}
  The average apparent velocity can be calculated by using the distance and
  time between the first and the last point which yields
  $\beta^{\mathrm{ave}}_{\mathrm{app}} = 4.8$.

\item[(T05.2a)] Radiation from the core always takes the same time to reach
  us, but the time taken by radiation from the moving blob in the jet
  decreases since its distance from us reduces.
  % FIGURE NEEDED: solution sketch — jet geometry triangle showing positions A
  % and B, viewing angle theta, and components of beta*c*Delta t along and
  % across the line of sight (2025 theory solutions, pages 13-14 of 59).

  As shown in the sketch above, a signal emitted from the moving blob at B
  covers a distance $(\beta c)\Delta t \cos\theta$ less than a signal emitted
  at A.

  Thus the apparent time difference between the light received by the
  observer is
  \begin{align*}
  \Delta t_{\mathrm{app}} = t_B - t_A
  &= \left(t_0 + \Delta t - \frac{(\beta c)\Delta t \cos\theta}{c}\right)
     - t_0 \\
  &= \Delta t\,(1 - \beta\cos\theta) \\
  \Longrightarrow\;
  \frac{\Delta t_{\mathrm{app}}}{\Delta t}
  &= (1 - \beta\cos\theta).
  \tag*{(2.0)}
  \end{align*}

\item[(T05.2b)] The apparent velocity on the sky is
  \begin{align*}
  \beta_{\mathrm{app}} = \frac{v_{\mathrm{app}}}{c}
  &= \frac{\beta\,\Delta t \sin\theta}{\Delta t_{\mathrm{app}}} \\
  \beta_{\mathrm{app}}
  &= \frac{\beta\sin\theta}{1 - \beta\cos\theta}
  \tag*{(2.0)}
  \end{align*}
  The resulting formula indicates that the measured apparent velocity of the
  jet's component is increased for relativistic features and can be higher
  than speed of light for small values of angle $\theta$.

\item[(T05.3a)] For $\beta_{\mathrm{app}} = 1$, we obtain
  \[
  \beta(\sin\theta + \cos\theta) = 1
  \;\Longrightarrow\;
  \beta = \frac{1}{\sin\theta + \cos\theta}
  \]
  Superluminal motion occurs if $\beta_{\mathrm{app}} > 1$; then
  \[
  \beta(\sin\theta + \cos\theta) > 1.
  \]
  This defines the superluminal region which is to be shaded.
  % FIGURE NEEDED: solution plot of beta = 1/(sin theta + cos theta) versus
  % theta (0-90 deg), symmetric about 45 deg with minimum 1/sqrt(2), with the
  % superluminal region above the curve shaded with slanted lines
  % (2025 theory solutions, page 15 of 59).

\item[(T05.3b)] It may be seen from graph drawn in (T05.3a) that the minimum
  value of $\beta$ for which superluminal motion is apparent occurs at
  $\theta = 45^\circ$, i.e., $\beta = 1/\sqrt{2}$.
  \begin{equation*}
  \beta_{\mathrm{low}} = 1/\sqrt{2} \approx 0.707,
  \tag*{(1.0)}
  \end{equation*}
  and the corresponding orientation angle,
  $\theta_{\mathrm{low}} = 45^\circ$. \hfill(1.0)

\item[(T05.4)] We know
  \[
  \beta_{\mathrm{app}} = \frac{\beta\sin\theta}{1 - \beta\cos\theta}
  \]
  Thus
  \[
  \beta\left(\frac{\sin\theta + \beta_{\mathrm{app}}\cos\theta}
                  {\beta_{\mathrm{app}}}\right) = 1
  \]
  Since $\beta \le 1$, we have
  \begin{equation*}
  \frac{\sin\theta + \beta_{\mathrm{app}}\cos\theta}{\beta_{\mathrm{app}}} > 1
  \tag*{(1.0)}
  \end{equation*}
  or
  \[
  \beta_{\mathrm{app}} \tan(\theta/2) < 1.
  \]
  Therefore, the maximum viewing angle for a jet with a given
  $\beta_{\mathrm{app}}$ is
  \begin{equation*}
  \theta_{\mathrm{max}}
  = 2\tan^{-1}\!\left(\frac{1}{\beta_{\mathrm{app}}}\right)
  \tag*{(1.0)}
  \end{equation*}

\item[(T05.5)] Time-scale $\delta t$ of any substantial variability from a
  source cannot be shorter than the light crossing time of the source.

  This then gives an upper limit to the size $\lesssim c\,\delta t$, where $c$
  is the speed of light. Therefore,
  \begin{equation*}
  c\,\delta t \gtrsim 5 R_{SC} = 5 \times \frac{2GM}{c^2}.
  \tag*{(1.0)}
  \end{equation*}
  where, $M$ is the mass of the central compact object. Rearranging,
  \[
  M \lesssim \frac{c^3\,\delta t}{10G}
  \]
  Thus, upper limit on the mass of central compact core of quasar
  \begin{align*}
  M_{c,\mathrm{max}} &= \frac{c^3\,\delta t}{10G}
  \tag*{(1.0)} \\
  &= \frac{(2.998 \times 10^{8})^3 \times 3600}
          {10 \times 6.674 \times 10^{-11}}~\mathrm{kg}
   = 7.311 \times 10^{7}\,M_\odot \\
  M_{c,\mathrm{max}} &= 7 \times 10^{7}\,M_\odot
  \tag*{(1.0)}
  \end{align*}
\end{description}
